\section{Color Transfer}

Дан Рудерман та інші спеціалісти. розробили кольоровий простір, який називається
lαβ, що мінімізує кореляцію між каналами для багатьох природних сцен. Даний простір
був зроблений на аналізі того як сприймає та обробляє вхідну інформацію зорова система
людини природні сцени. Автори представили новий формат lαβ беручи за основу
зорову систему людини.

Між осями у просторі lαβ є вирази, що дозволяють нам використовувати різні дії в
різних каналах, при чому не переживаючи, що не з’являться дефекти на зображенні.
Плюс - цей колірний простір логарифмічний, що значить нормальне масштабування.

Даний алгоритм досі не розглядався у контексті роботи із супутниковими знімками.
Тому є необхідність зробити певний аналіз щодо того чи слід його використовувати
між сусідніми регіонами, оскільки до цього його застосовували тільки для звичайних
RGB зображень.

Було запрограмовано вищеописаний алгоритм та виявлено такі факти:

\begin{enumerate}
	\item \textbf{Не було знайдено методу, що переводить Geotiff знімок у формат LAB}
		Проблема полягає в тому, що оригінальний алгоритм перенесення кольору загалом
		працює із зображеннями з малим масштабом, а саме наприклад значення [0 - 255].
		У знімках формату Geotiff діапазон набагато більший [0 - 10 000];

	\item \textbf{Метод не може забезпечити правильну відповідність територій}
        Тут мається на увазі, що якщо на зображенні на яке ми хочемо накласти новий колірний
		простір є регіон, наприклад лісів чи полів, то після обробки ті ж самі
		регіони мають і залишитись лісами чи полем..
\end{enumerate}

Провівши аналіз алгоритму перенесення колірного простору, було прийнято рішення щодо
застосування методів машинного навчання, адже саме завдяки ним можна одержати
правильне зображення де у відповідності зі знімком з якого беруть колір на
знімку до якого його застосували будуть максимально наближені регіони по своєї
притаманності, тобто якщо було озеленене поле воно стане більш сухим і світлим,
якщо такі були наявні у знімку з якого брали стиль.

Коли мова йде про обробку супутникових зображень і щоб на виході отримувати не
класифікатор, а значення пікселів, то загалом використовують регресійні методи
машинного навчання.

До найбільш придатних алгоритмів які вміють працювати зі знімками відносять
лінійну, поліноміальну, баєсову регресії, регресію в методі опорних векторів та відомий
багатошаровий перцептрон і алгоритм випадкового лісу. Вважаючи доцільним їх
застосування у даній роботі, потрібно зробити їхній огляд.
